\clearpage\SourceLeaf{185}
\begin{center}[ASTRALGEOMETRIE.]\end{center}

\begin{center}\rule{3cm}{0.4pt}\end{center}

\begin{center}[\textsc{Gerling} \emph{an} \textsc{Gauss}. \emph{Marburg, 23. Juli 1818.}]\end{center}

\{.\ .\ .\ .\ .\ .\ Im nächsten Semester steht mir eine Arbeit bevor, für die ich
mir schon jetzt erlaube, Sie um Rat zu bitten. Wahrscheinlich werde ich mir
das noch öfter erlauben. --- \textsc{Fleckeisen} hat mich
nämlich gebeten, eine neue Auflage von \textsc{Lorenz}’
reiner Mathematik zu betreuen. Ich tue das umso lieber, weil
ich hier selbst nach diesem Buch lehre und es nun schon 6 Jahre lang
im Unterricht verwende.

.\ .\ .\ .\ .\ .\ .\ In der Geometrie habe ich weniger Einwände, würde aber besonders
gern Ihre Meinung dazu hören, wie man mit der Parallelentheorie am besten
umgeht. Was \textsc{Lorenz} dazu schreibt, ist teils falsch, teils
unzureichend begründet. --- Ich halte es für richtig, die euklidische Vorgehensweise darzustellen und dabei
ihren Mangel einzugestehen. Das quomodo ist
mir aber nicht klar. Ich dachte, es sei am besten, den Satz, dass eine gerade
Linie durch einen Punkt nur eine Parallele haben kann, als Axiom
aufzustellen und in einer Anmerkung zu sagen, dass man für
diesen Satz noch keinen Beweis finden konnte und ihn deshalb, bis
ein Beweis gefunden oder der Satz widerlegt ist, als Axiom
annehmen muss, wie es im Grunde schon Euklid getan hat. --- Bitte haben Sie die
Güte, mir auch hierzu Ihr Urteil mitzuteilen. .\ .\ .\ .\ .\ .\}

\begin{center}\rule{3cm}{0.4pt}\end{center}

